\documentclass[conference]{IEEEtran}
% Independent R1 revision. Based on the audited supplied PDF and SE25 source;
% neither is asserted to be the exact compilation source of the other.
% r1-results.tex and figures are generated by the registered experiment workflow.
\usepackage[T1]{fontenc}
\usepackage[utf8]{inputenc}
\usepackage{cite,amsmath,amssymb,graphicx,booktabs,array,url}
\usepackage[hidelinks]{hyperref}
\newcommand{\op}[1]{\texttt{#1}}
\newcommand{\RoneTestCount}{423}
\newcommand{\optionalfigure}[3]{\IfFileExists{figures/#1}{%
  \begin{figure}[tb]\centering\includegraphics[width=\columnwidth]{figures/#1}%
  \caption{#2}\label{#3}\end{figure}}{}}
\hypersetup{pdftitle={Predictive Multi-Contact Security Lifecycle Management: R1 Evidence Revision},pdfauthor={Mukhammed Zhumali}}
\begin{document}
\title{Predictive Multi-Contact Security Lifecycle Management for LEO Spacecraft--Ground Communications}
\author{\IEEEauthorblockN{Mukhammed Zhumali}\IEEEauthorblockA{262937@astanait.edu.kz}}
\maketitle

\begin{abstract}
Time-limited spacecraft--ground contacts constrain payload transmission and the operations needed to maintain traffic keys and security associations. We study when additional contact visibility changes lifecycle decisions and what management and computational costs it incurs. Matched Next and Multi policies share a finite action model, executor, constraints, five-component objective and full-plan tie rule; only the permitted horizon differs. A registered campaign produces 135 finished executions in 181 attempts, with 55 of 84 selected horizon pairs complete. A completed exact short-return witness gains one served contact with Multi at higher management cost and one second less useful transmission. Bounded practical search retains a gain in two finite-demand cases but loses that gain under saturated demand and can perform worse on a high-demand fragment. Full continuous ISS, NOAA-19 and HST search-policy runs exceed the registered budgets. Separate interventions examine decision delay, traffic-key disclosure and actual channel interruption; aborted committed frames consume their nonces. The model assumes a logical ground coordinator, synthetic operation costs and protected management credentials. It supplies neither a complete distributed protocol nor onboard validation. The results support a conditional mechanism and identify computational limits, rather than universal policy superiority or improved cryptographic strength.
\end{abstract}
\begin{IEEEkeywords}
LEO communications, key lifecycle, security association, contact-aware planning, bounded search, disruption, receding horizon
\end{IEEEkeywords}

\section{Introduction}
A return contact can be long enough to transmit its required payload but too short to install, activate, and use a replacement key. If that requirement is visible during an earlier contact with the same station, some work can be moved forward. The scheduling opportunity depends on key validity, resource contention, payload requirements, and the selected plan. A longer forecast alone guarantees neither better service nor lower cost.

The research question is therefore conditional: \emph{when does additional contact visibility change joint traffic-key and Security Association (SA) decisions, and what service benefit, management cost, and computational burden follow?} Multi-contact denotes observation depth for one spacecraft and several ground stations. It does not denote a multi-spacecraft optimization problem.

The contribution is a finite lifecycle-planning formulation and an executable comparison that separates forecast depth from other policy choices. The primary objective is admissible service, not compromise probability. Disclosure experiments examine consequences of chosen schedules; they do not introduce advance knowledge of an attack or optimize recovery latency. Receding-horizon optimization, contact prediction, advance authentication, and established cryptographic mechanisms are not claimed inventions.

R1 addresses two applicability gaps. First, it separates an exact finite-model reference from a practical policy that may use a feasible but uncertified incumbent or a common fallback. Second, it makes decision delay and actual channel unavailability explicit. These extensions test the price and limitations of the mechanism. They do not turn synthetic timing assumptions into measurements of a flight implementation. Detailed historical version and retry information is placed in the separate supplement.

\section{Related Work and Novelty Boundary}
CCSDS SDLS specifies data-link protection, and its Extended Procedures define key and SA management services~\cite{sdls,ep}. The Symmetric Key Management Recommended Practice supplies an explicit key lifecycle and abstract management procedures~\cite{skm}. Lifecycle states, key distribution, activation, and retirement therefore pre-exist this study. We schedule abstract counterparts under a chosen resource and service model; sharing state names does not establish SDLS conformance.

Connectivity already informs security-state management. The IOTA Identity architecture updates an onboard authorization allowlist during ground contacts and considers revocation delay and short-lived capability tokens~\cite{iota}. Its object is IoT-device identity and access authorization. Our object is directional traffic-key and SA timing. The source audit also retains FHAP and LK-AKA as authentication precedents, but publisher access to FHAP and full-text access to LK-AKA were limited. We do not infer detailed scheduling or experimental omissions from unavailable full texts.

Delay-Tolerant Networking supplies a second relevant line of work. BPv7 provides bundle delivery in challenged networks~\cite{bp}; BPSec provides bundle integrity/confidentiality while leaving key management to separate network management~\cite{bpsec}. Schedule-Aware Bundle Routing explicitly uses scheduled contacts for route computation~\cite{sabr}. Hence prediction and contact-based optimization alone are insufficient novelty claims.

Most directly, Detke et al. couple time-feasible Contact Graph Routing routes with BPSec key provisioning~\cite{detke}. Their offline mixed-integer program selects routes, active key scopes, and covered pairs to minimize scopes at a target connectivity ratio; temporal routing feasibility is an input, while the displayed variables do not schedule key/SA operations or encode cryptoperiods. Their evaluation uses 20 nodes, four domains, and 64 contacts over 2700~s, and assesses provisioning, processing, and data-at-rest exposure. They explicitly identify key scheduling over contact plans as future work. Our distinction is sequential per-station lifecycle action timing and controlled visible-horizon comparison for a fixed workload, rather than priority over all temporal DTN key-management formulations.

\begin{table*}[t]
\centering\small
\caption{Closest established mechanisms and the bounded comparison made here. No external protocol was reimplemented or experimentally outperformed.}
\label{tab:related}
\begin{tabular}{p{.18\textwidth}p{.32\textwidth}p{.43\textwidth}}
\toprule
Source & Established object/mechanism & Relationship to this study \\
\midrule
SDLS/EP and symmetric key management & Space-link protection; key lifecycle; concrete and abstract management procedures & Supply mechanism/state provenance. R1 studies timing under explicit costs and constraints, not standards compliance. \\
IOTA Identity & Device credentials, contact-dependent allowlists, revocation handling & Establishes connectivity-aware authorization. R1 concerns channel traffic keys and paired SAs. \\
BPv7, BPSec, SABR & Bundles, bundle security, and scheduled-contact routing & Establish intermittent-connectivity processing and prediction-based routing. R1 fixes the selected link workload. \\
Detke et al., full-paper comparison & Time-feasible route catalog coupled to offline BPSec scope provisioning & Closest reviewed optimization precedent. R1 varies observation depth for sequential lifecycle execution; Detke leaves contact-plan key scheduling to future work. \\
\bottomrule
\end{tabular}
\end{table*}

\section{System and Lifecycle Model}
\subsection{Security State and Validity}
One spacecraft maintains an authorized context for each configured ground station. Uplink and downlink have separate traffic keys and SAs; the two SAs are managed as a pair. At most one ground-link context is operational. Station visibility is not authorization, and keys are not silently reused across station contexts.

The implementation stores immutable metadata and state, not secret key bytes or cipher execution. Preparation, confirmed installation, activation/binding, and opening a session are distinct. A prepared key does not authenticate a peer. Closing a session preserves valid keyed state and persistent counter progress. Retired identities are not recycled. Logical retirement is not evidence of physical erasure.

For creation time $t_k^c$, first activation $t_k^a$, material-age bound $A_{\max}$, activation cryptoperiod $C_{\max}$, and committed not-after bound $\bar d_k$, the active key deadline is
\begin{equation}
d_k=\min\{\bar d_k,t_k^c+A_{\max},t_k^a+C_{\max}\}.
\end{equation}
The activation term does not apply before activation. Reopening or reusing an SA does not extend an existing deadline. Between-contact reuse also requires that the gap from the last successful directional session does not exceed $G_{\max}$, together with authorization, installation, peer synchronization, and counter headroom.

Online confirmations must finish strictly before contact closure and required-key expiry. A complete payload frame may end exactly at either boundary. Expiry is an actual execution event. If one directional key expires, the common response stops the paired context and retires its companion according to the model. Independent staged material is retained when admissible.

\subsection{Execution and Costs}
The executor supports \op{PREPARE}, \op{INSTALL}, \op{SWITCH}, \op{OPEN}, \op{REUSE}, one-frame \op{CONTINUE}, and \op{WAIT}. STOP and expiry are local transitions; incident response is common to all policies. Operations reserve ground CPU, spacecraft CPU, radio direction, and waiting phases. Only executed work is charged, including partial or wasted work. Useful time and delivered service count complete frames.

\begin{table}[t]
\centering\small
\caption{Base synthetic operation profile; not hardware calibration.}
\label{tab:costs}
\setlength{\tabcolsep}{4pt}
\begin{tabular}{lrrrrr}
\toprule
 & Elapsed & GS CPU & SC CPU & UL & DL \\
 & s & s & s & B & B \\
\midrule
PREPARE & 2.0 & .20 & .00 & 0 & 0 \\
INSTALL & 4.0 & .10 & .20 & 600 & 200 \\
SWITCH & 1.0 & .05 & .10 & 100 & 100 \\
OPEN/REUSE & 1.0 & .02 & .04 & 40 & 40 \\
Revocation notice & .4 & .02 & .02 & 64 & 32 \\
\bottomrule
\end{tabular}
\end{table}

The profile in Table~\ref{tab:costs} models a 400~B/s management channel. Elapsed duration includes waiting and is not a sum of CPU measurements. PREPARE uses ground resources only and may overlap payload if reservations permit; online management competes for the link. STOP and expiry have zero modeled cost in this profile. These values are controlled assumptions, not values prescribed by CCSDS or measured on a spacecraft.

The state includes pending operations, committed frames, per-contact delivery, and accrued costs. Replanning does not cancel committed work or charge it again. Confirmed state changes occur only at operation checkpoints. An uncertain SWITCH, OPEN, or REUSE is blocked pending explicit reconciliation; restoration of connectivity alone is not reconciliation.

\subsection{Deployment Interpretation}
R1 assumes a \emph{ground logical coordinator} with access to the permitted forecast and a consistent abstract execution state. Ground-only preparation can occur while disconnected. Installation, switching, opening, reuse, and spacecraft notice require the matching available link. Local stop and expiry do not require successful delivery of a new command.

This specifies placement and the timing boundary but closes deployment only partially. It assumes trusted management credentials and abstract confirmations. It does not implement distributed telemetry/state reconciliation, an onboard command protocol, failover between controllers, or measured command/telemetry bytes. The single logical state is an explicit abstraction, especially when a downlink continues after ground revocation notice. A flight realization must map each abstract operation to a concrete protocol and validate its costs and information flow.

\section{Horizon, Objective, and Policy Modes}
\subsection{Observation and Objective}
For a decision in contact $C_i$, $H$ counts \emph{additional future contacts}:
\begin{equation}
\mathcal W_i^H=(C_i,C_{i+1},\ldots,C_{\min(n,i+H)}).
\end{equation}
Thus Next uses $H=1$ and sees current plus next; the basic Multi comparison uses $H=2$ and sees current plus two. In a gap there is no current contact, so the nearest $H$ future windows are visible. A search freezes this observation and its endpoint; advancing a hypothetical branch does not reveal an additional contact.

Future prescribed actions, channel outages, restoration times, attack occurrences, and unobserved failures are excluded. Nominal branches model future confirmations as successful until actual execution supplies contrary evidence. Only the next selected action is admitted before the controller observes and replans.

Under the common admission constraints, the objective is
\begin{equation}
\max_{\mathrm{lex}} J=(N_{\mathrm{served}},T_{\mathrm{useful}},-C_{\mathrm{CPU}},-B_{\mathrm{mgmt}},-N_{\mathrm{new}}).
\label{eq:objective}
\end{equation}
Service requires the configured complete-frame minimum separately in each direction. CPU and management bytes sum both endpoints/directions; new keys count directional identities, including wasted preparation. Previously accrued work is retained and already served contacts are not counted twice. There is no terminal bonus for unused fresh material. Recovery latency, readiness, and compromise probability are not additional objectives.

Only equality of all five components invokes a full-plan tie rule. CANONICAL orders preparation, installation, switching, opening, reuse, continuation, and wait, then start time and the full identity/end-time tuple. Plans are compared at their first difference, with a proper prefix preceding its extension. CONTINUE\_WAIT places continuation, then wait, then management first; within a group it compares earlier end time before the original tuple. It is more specific than an action-order permutation. Search traversal order is distinct from this preference. Within each rule and execution mode, Next and Multi share all mechanisms except horizon.

\subsection{Exact Reference and Practical Execution}
The exact reference enumerates a finite event-based action model. Candidate start opportunities are discrete; at most one spare pair or preparation is admitted per context; frame direction follows the common executor; actions must satisfy completion-within-horizon restrictions. Fully solved roots are exact only within this model. Neither these root solutions nor cached equivalent suffixes prove global optimality of the executed receding-horizon policy.

A root status COMPLETE certifies its finite continuation; INCOMPLETE does not. The reference controller stops on an incomplete search and preserves a checkpoint instead of executing an unproved fallback. FINISHED denotes a completed full scenario and is a different status. The practical policy is deliberately separate: for a local node/time limit it may execute a feasible terminal incumbent's first action, or a common admissible service/bootstrap fallback followed by WAIT when necessary. Its original incomplete search and disposition remain recorded. Cancellation or an outer controller stop does not become permission to continue.

The fallback is a specified execution rule, not a sixth objective or a certificate of optimality. A practical FINISHED result may contain incomplete searches; report both completion and the fraction/type of incumbent/fallback decisions. A practical/exact difference is a policy-mode comparison, not an isolated forecast-depth effect.

Practical-v2 may retain the feasible incumbent after one completed frame strictly inside the same contact. Its next observed state must equal the entire nominal observation, including visible contacts, horizon, costs, counters, progress, issued identifiers and committed work; the exact next action must still be admissible. A service action ends this retention chain. The source INCOMPLETE result remains in the audit, and no optimality certificate is created. This differs from certified suffix reuse after a completed exact search.

\begin{table}[t]
\centering\small
\caption{Registered search controls. Host limits are separate from simulated operation costs.}
\label{tab:budgets}
\begin{tabular}{lrr}
\toprule
Control & Exact reference & Practical \\
\midrule
Nodes per search & 2,000,000 & 5,000 \\
Elapsed per search & 30 s & 1 s \\
Controller elapsed & 60 s & 20 s \\
Nominal process watchdog & 90 s & 35 s \\
\bottomrule
\end{tabular}
\end{table}

\subsection{Operational Baselines}
FixedPeriod is an author-defined operational baseline, not a mandated CCSDS scheduler. Its renewal timer is anchored to the last successful SWITCH; overdue requests coalesce, and a missed renewal does not extend key validity. Historical X/Y and SE2 ISS settings use 60 and 1800~s, respectively. Earlier ISS controls used 3600~s and retain separate provenance. Historical sensitivity also includes a periodic setting that serves every X contact. One selected period cannot establish superiority over all periodic policies.

PerContact requires newly activated material for the contact under its registered preparation rule. Both operational controls share admission checks and the event executor. They provide context; the primary horizon attribution remains the matched Next/Multi comparison.

\section{Delay and Channel-Availability Experiments}
\subsection{Decision and Dispatch Delay}
Offline reference search does not advance physical simulation time. R1 separately models a critical-path delay
\begin{equation}
\delta=\delta_{\mathrm{fixed}}+\alpha t_{\mathrm{host}},\qquad \delta_{\mathrm{fixed}},\alpha\geq0,
\end{equation}
where $t_{\mathrm{host}}$ is measured selection time and the other terms are experimental assumptions. The fixed term includes the chosen decision/dispatch allowance; it is not a measured radio protocol delay. The coordinator issues no new action while computing, while already committed work and physical events continue.

After the delay, the controller constructs a fresh observation and revalidates the first selected intent. A deadline miss or unavailable/stale intent follows the registered WAIT behavior rather than applying a stale optimal plan. An observed incident supersedes the scheduled choice through the common response. This shifted execution has no inherited optimality certificate. Aggregate host runtime cannot simply be subtracted from a service or recovery benefit; only the relevant critical-path delay has that interpretation. No onboard CPU equivalence is inferred from $\alpha$.

\subsection{Controlled Outages Without Future Leakage}
An interruption names a contact and an actual unavailable interval wholly inside its original forecast window. Full-window loss represents a missed contact; initial loss represents delayed availability; trailing loss represents early closure. The original forecast remains unchanged.

Actual DOWN aborts a committed frame, fails an online primitive with its consumed cost, and stops operational use locally. PREPARE is ground-only and remains possible. Actual UP records availability but does not automatically reopen an SA, confirm an interrupted step, or clear uncertainty. INSTALL can be retried if subsequently admissible; uncertain SWITCH/OPEN/REUSE remains fail-closed because reconciliation is not implemented.

The observation first selects the original H-limited prefix, then removes a currently unavailable contact based only on past DOWN/UP events. It does not fetch another farther contact to replace it. The branch receives no future restoration event. Actual UP wakes execution and allows a new observation. This is conservative handling of known loss, not estimation of its duration or a complete model of forecast error.

Frame completion precedes loss at the same timestamp; expiry and original contact closure precede online confirmation. Loss precedes an ACK at the same instant. UP precedes DOWN for adjacent outage intervals, avoiding a fictitious available instant. These conventions are explicit simulation choices.

\subsection{Aborted-Frame Accounting Correction}
The earlier abort path recorded emitted ciphertext but did not consume the committed directional nonce. R1 corrects this before allowing a temporarily interrupted context to reopen with the same key. Each aborted committed frame now increments its key counter once and updates any still-bound SA mirror, while contributing zero delivered frames, bytes, and useful time. Even a zero-byte prefix conservatively consumes that nonce.

The operation is accounting for an already admitted transmission, not new permission to transmit. A ground-rejected committed downlink can consume its retired key's counter without resurrecting that key or refreshing successful-use history. New counter audits therefore count delivered \emph{and aborted committed} frames. Historical files are preserved and must retain their executable identity; this correction is not retroactively claimed for old runs.

\section{Experimental Design and Metrics}
\subsection{Workloads and Registered Comparisons}
The controlled X witness has contacts A--B--A, with a short return after the initial A generation expires. Y changes the return station. Historical X return durations of 3/5/9~s isolate remaining-work feasibility, not an independently sampled mission distribution. Geometric workloads use public orbital elements and station locations with a common selector independent of security policy.

The R1 preparation adds public GP inputs for NOAA~19 and HST to broaden geometric workload checks. The obtained base schedules contain 38 and 48 selected windows, respectively. Geometry generation is not a completed policy run or evidence of authorized communications. Full execution coverage, selected fragments, hashes, and completion are reported only in the generated result block. Actual radio-link conditions, station availability, and authorization are not inferred from geometry.

Within a paired comparison, initial state, workload, traffic, operation costs, validity/reuse constraints, observed incidents, execution mode, tie rule, and budgets are shared; only the registered factor changes. Exact/practical, latency, and outage interventions are identified separately. Search budgets and failed/incomplete attempts are retained. A duration, cost, load, or horizon sweep must report its evaluated grid and missing cells rather than select only favorable cases.

\subsection{Service, Security, and Computation}
Alongside the five objective components, report full execution status, search status, wall time, explored nodes, and practical incumbent/fallback use. Computational effort is not the modeled security-operation CPU cost. A process failure, search limit, and service failure are distinct.

Disclosure experiments reveal a selected active traffic-key pair to a recording attacker. Replacement material is assumed independent; management credentials remain uncompromised. A notice is exogenous. Recovery is separated into spacecraft notice, independent-pair activation, independent-SA opening, and first independent delivered frame. Missing milestones in completed executions are right-censored, not zero. Incomplete executions have no final quality result.

Confidential service is retrospective: each direction's delivered minimum must be met by complete frames whose keys remain undisclosed at experiment end. This excludes earlier traffic under a subsequently disclosed key and never feeds future knowledge into planning. It does not mean every byte in the contact is confidential. Emitted exposure includes aborted prefixes, so exposure and confidential delivery are not a partition of a common delivered-byte total.

\section{Results and Evidence Status}
\subsection{Historical Evidence, Separate from R1}
The saved October 24 synthesis contains 84 selected executions: 63 FINISHED, 20 search-incomplete, and one controller-limited. Only 21 of 42 matched pairs finish both policies; three further technical attempts are stored separately. Twenty X/5 rows are retained historical runs, not new independent trials. A current CSV recount agrees with these counts; that recount is not a fresh simulation or revalidation of every underlying journal.

The autonomous nominal X/5 witness serves two contacts with Next and three with Multi, with 38 versus 37~s useful transmission and higher management cost for Multi. The service difference survives both tested full-plan selectors. Completed X/9 CONTINUE\_WAIT comparisons become equal, while other unfinished cells remain unknown. In the two additional ISS fragments, canonical early installation does not improve served contacts or confidential volume. Equal-score choices affect readiness and recovery timing. These findings support a conditional mechanism, not universal resource savings, recovery superiority, or reduced compromise probability.

For two specified return-entry states, the mechanism can be expressed as
\begin{align}
\tau_{\mathrm{installed}}&=t_S+t_O+t_P,\\
\tau_{\mathrm{prepared}}&=t_I+\tau_{\mathrm{installed}}.
\end{align}
Here $t_I,t_S,t_O$ are installation, switching, and opening durations, $t_P$ is the time for the required complete-frame payload, and $W$ is remaining contact time. When $\tau_{\mathrm{installed}}\leq W<\tau_{\mathrm{prepared}}$, only the installed-entry chain fits. This is a feasibility statement for those chains, requiring all validity/resource checks. It is not a lower bound over all possible policies or initial states.

\subsection{R1 Experiments}
\IfFileExists{r1-results.tex}{\subsubsection{Coverage and execution mode}

R1 attempted 181 registered post-change executions; 135 reached FINISHED. Of 84 matched horizon pairs, 55 finished both policies. The 12 pre-change baseline attempts produced 11 FINISHED and one search-incomplete execution and are excluded from these R1 counts.

\begin{table}[t]\centering\caption{R1 execution completion. Search completeness is separate.}\label{tab:r1coverage}\begin{tabular}{lrrr}\toprule Series & Planned & Attempted & Finished\\\midrule

exact & 40 & 40 & 30\\

practical & 44 & 44 & 23\\

security & 24 & 24 & 24\\

disruptions & 16 & 16 & 16\\

sensitivity & 38 & 38 & 26\\

controls & 10 & 10 & 10\\

deadlines & 8 & 1 & 0\\

deadlines-v3 & 8 & 8 & 6\\

\bottomrule\end{tabular}\end{table}

\subsubsection{Finite-budget horizon effects}

The exact series finishes 30 of 40 executions and 10 of 20 pairs. Eight searches stop incomplete and two controllers reach their total budget. The finished exact X/5 CONTINUE\_WAIT pair retains three versus two served contacts for Multi versus Next, at 37 versus 38 s useful time, 0.83 versus 0.12 s operation CPU, 1240 versus 160 management bytes, and two versus zero new keys. This is the autonomous preparation model, not the older shared-preparation accounting.

At the practical 5000-node/1-s decision budget, both rules produce only two served contacts for both policies on X/5: the earlier exact benefit is lost. On the high-demand no-attack ISS fragment, CANONICAL serves three contacts for both policies, but Multi spends 0.65 s more operation CPU, 1000 additional management bytes and creates two additional keys. CONTINUE\_WAIT serves two contacts with Multi versus three with Next (40 versus 410 s useful time). These are feasible bounded-controller outcomes, not a reversal of an exact optimum; additional horizon competes for the same search budget.

Finite offered demand changes this outcome: at 400 and 1600 bytes per direction/contact, CANONICAL practical Multi gains one served contact, with respectively 2 and 4 s extra useful time. Each gain costs 0.71 s more operation CPU, 1080 additional management bytes and two new keys. Those factor cells do not test the other tie rule. FixedPeriod T=90 s also serves all three X/5 contacts; all tested full-ISS operational controls serve 46 contacts. Full search-policy execution remains incomplete, so no missing final service score is ranked against those controls.

Tightening cryptoperiod to 90 s forces the old A pair to expire within C1 at t=10; Next activates its replacement at t=15 and both horizons then serve three contacts. This is a case where current-contact renewal makes extra visibility unnecessary. Increasing cryptoperiod to 240 s does not extend a stored earlier notAfter. These specific outcomes do not imply monotonic benefit or arbitrary-cryptoperiod robustness.

\subsubsection{Diagnostics, controls and model boundaries}

A fixed-root 20000-node diagnostic preserves the entire incumbent plan and all five metrics under both rules, but the aggregate change is 1.89\% and 2.47\% slower in these individual profiled runs. No acceleration is claimed. Twelve practical-versus-full-DFS comparisons from two previously established warmed states have zero difference in all objective components; the initial fresh-root oracle exceeded its unchanged node cap. This local agreement supplies no general approximation guarantee.

The matrix separately retains undetected and delayed disclosures, gap disclosure, four channel interventions, fixed-period controls, and synthetic 50/500 ms decision delays against a 250 ms deadline. Journal-derived recovery distinguishes spacecraft notice, independent activation, SA opening and first delivery. Successful activation does not imply continuous readiness. Full-sequence ISS, NOAA-19 and HST attempts retain incomplete outcomes; they are not replaced with selected triples. Numeric sensitivity and per-case source links are available in the reproduction tables.
}{%
The registered R1 result fragment has not yet been integrated into this manuscript. No new matrix outcome, robustness gain, runtime advantage, or practical-policy service result is asserted here. Geometry preparation and implementation checks are distinct from completed policy experiments.}

\optionalfigure{runtime_completion.pdf}{R1 computational effort and completion coverage. Incomplete executions retain their status; any runtime summaries must state their completion conditioning.}{fig:runtime}
\optionalfigure{service_cost_horizon.pdf}{Service and resource cost by permitted horizon for the registered R1 comparisons. Interpretation is restricted to the displayed configurations and completion coverage.}{fig:horizon}
\optionalfigure{sensitivity.pdf}{Registered sensitivity outcomes. Missing or incomplete cells are retained and are not assigned policy-quality values.}{fig:sensitivity}
\optionalfigure{mechanism_timeline.pdf}{Lifecycle timing for the identified workload and policy pair. Operation timing explains a mechanism; it does not alone establish a general performance claim.}{fig:mechanism}
\optionalfigure{recovery.pdf}{Separate recovery milestones under the registered disclosure assumptions. Censored and incomplete outcomes require distinct markings.}{fig:recovery}
\optionalfigure{detected_exposure.pdf}{Detected disclosure in the CANONICAL practical X cases. Colors identify policies; hatching identifies post-disclosure exposure, with recorded prior ciphertext shown separately. Equal bars are retained.}{fig:exposure}

\subsection{Implementation Verification}
At the integration checkpoint, the fresh candidate suite records \RoneTestCount{} passing tests. This verifies the exercised implementation and fixtures, not statistical generalization, a cryptographic proof, or all ISS-sized optima. Added checks cover observation noninterference, interrupted management work, counter accounting, conservative uncertainty, equal-time boundaries, and checkpoint restoration. Final executable identity and any subsequent integration checks belong to the experiment manifest and supplement; test totals from historical stages are not added to this suite count.

After practical-v2, 74 overlapping focused Java tests and 67 Python tests pass. A subsequent export repair passes 10 focused tests. They include exact-reference comparisons on two warmed witnesses, bounded feasible retention, cancellation, outage restoration, and corrupt-journal rejection. These counts are not added to the earlier suite total.

\section{Discussion and Limitations}
The benefit of visibility is mediated by available earlier contacts, remaining material validity, whole-frame demand, and management-resource contention. Installing early can preserve later service while reducing present payload or increasing cost. When deferred work fits the return window, or offered demand is already satisfied, earlier readiness can produce no delivered-service gain. A nominally equal alternative plan can also change incident timing without changing Eq.~\eqref{eq:objective}.

The exact reference is an offline finite-model benchmark. The practical mode trades certification for bounded decision effort and therefore must be evaluated as an explicitly different controller. Its completed executions do not imply solved roots. Delay sensitivity exposes selected critical-path assumptions but does not measure a complete ground/onboard deployment. Controlled losses test specific interventions; they do not identify a general distribution of contact errors or detector behavior.

Operation and traffic costs remain synthetic. Public propagated visibility provides geometric realism but not operational access, link budgets, authorized ISS/NOAA/HST sessions, or independent mission samples. Cold-start fragments and deliberate return-pattern selection limit mission-level inference. High-demand or large-horizon incomplete executions cannot support a claim that Multi is better, equal, or never worse.

The threat model excludes persistent generator compromise, active forgery, a physical RF jammer model, compromised management credentials, and a real distributed endpoint protocol. The nonce correction closes one concrete accounting gap; it does not implement encryption or prove replay resistance. Explicit reconciliation after uncertain online operations remains absent. These are limitations of the current evidence and architecture, not completed features.

\section{Conclusion}
Additional contact visibility can reveal a renewal requirement early enough to move installation into a preceding contact. Completed exact and finite-demand practical cases support that conditional mechanism and its resource price. Saturated practical cases, high-demand losses and incomplete continuous workloads limit its applicability under the tested budgets. Shorter validity can instead force Next to renew within the current contact, removing the need for additional visibility. R1 separates finite-model exactness, feasible execution, decision delay and actual channel availability while preserving the objective and information boundary. Further work must address measured search costs and distributed reconciliation, with calibrated mission parameters. The study establishes neither universal superiority nor deployment readiness.

\begin{thebibliography}{9}
\bibitem{sdls} CCSDS, \emph{Space Data Link Security Protocol}, CCSDS 355.0-B-2, July 2022. [Online]. Available: \url{https://ccsds.org/Pubs/355x0b2.pdf}.
\bibitem{ep} CCSDS, \emph{Space Data Link Security Protocol---Extended Procedures}, CCSDS 355.1-B-1, Feb. 2020. [Online]. Available: \url{https://ccsds.org/Pubs/355x1b1.pdf}.
\bibitem{skm} CCSDS, \emph{Symmetric Key Management}, CCSDS 354.0-M-1, Dec. 2023. [Online]. Available: \url{https://ccsds.org/Pubs/354x0m1.pdf}.
\bibitem{iota} D. Bernal, O. Ledesma, P. Lamo, and J. Bermejo, ``IOTA-Based Authentication for IoT Devices in Satellite Networks,'' \emph{Computers, Materials \& Continua}, vol. 86, no. 1, pp. 1--39, 2026, doi: 10.32604/cmc.2025.069746.
\bibitem{bp} S. Burleigh, K. Fall, and E. Birrane III, ``Bundle Protocol Version 7,'' RFC 9171, Jan. 2022. [Online]. Available: \url{https://www.rfc-editor.org/rfc/rfc9171.html}.
\bibitem{bpsec} E. Birrane III and K. McKeever, ``Bundle Protocol Security (BPSec),'' RFC 9172, Jan. 2022. [Online]. Available: \url{https://www.rfc-editor.org/rfc/rfc9172.html}.
\bibitem{sabr} CCSDS, \emph{Schedule-Aware Bundle Routing}, CCSDS 734.3-B-1, July 2019. [Online]. Available: \url{https://ccsds.org/Pubs/734x3b1.pdf}.
\bibitem{detke} R. Detke, P. G. Madoery, R. Cherini, J. A. Fraire, and J. M. Finochietto, ``Security-Aware Route Selection in DTNs: Towards Optimal Key Provisioning,'' IEEE WiSEE, Space-Terrestrial Internetworking Workshop (STINT), 2026. Conference-hosted full paper, accessed Oct. 25, 2026. [Online]. Available: \url{https://stintworkshops.org/papers/stint2026-19-detke.pdf}.
\end{thebibliography}
\end{document}
